% !TEX program = pdflatex
\documentclass[11pt, a4paper]{book}

\usepackage[T1]{fontenc}
\usepackage{palatino}
\usepackage{mathrsfs}
\usepackage{amsmath, amssymb, amsthm}
\usepackage{geometry}
\usepackage{microtype}
\usepackage{setspace}
\usepackage{titlesec}
\usepackage{hyperref}
\usepackage{mdframed}
\usepackage{longtable}
\usepackage{booktabs}
\usepackage{array}
\usepackage{enumitem}
\usepackage{xcolor}
\usepackage{tocloft}
\usepackage{fancyhdr}

\geometry{top=1.1in, bottom=1.1in, left=1.25in, right=1.25in}
\onehalfspacing

\hypersetup{colorlinks=true, linkcolor=black, citecolor=black, urlcolor=black}

% Theorem environments
\newtheoremstyle{flyxion}
  {10pt}{10pt}{\itshape}{}{\bfseries}{.}{.5em}{}
\theoremstyle{flyxion}
\newtheorem{theorem}{Theorem}[chapter]
\newtheorem{proposition}[theorem]{Proposition}
\newtheorem{lemma}[theorem]{Lemma}
\newtheorem{corollary}[theorem]{Corollary}
\newtheorem{conjecture}[theorem]{Conjecture}
\newtheorem{definition}[theorem]{Definition}
\newtheorem{claim}[theorem]{Claim}

% Remark style
\theoremstyle{remark}
\newtheorem{remark}[theorem]{Remark}
\newtheorem{openq}[theorem]{Open Question}

% Box environments
\newmdenv[linewidth=0.8pt, innertopmargin=10pt, innerbottommargin=10pt,
  innerleftmargin=13pt, innerrightmargin=13pt,
  skipabove=12pt, skipbelow=12pt]{flybox}

\newmdenv[linewidth=1.2pt, innertopmargin=11pt, innerbottommargin=11pt,
  innerleftmargin=14pt, innerrightmargin=14pt,
  skipabove=12pt, skipbelow=12pt, backgroundcolor=gray!5]{masterbox}

\newmdenv[linewidth=0.6pt, innertopmargin=9pt, innerbottommargin=9pt,
  innerleftmargin=12pt, innerrightmargin=12pt,
  skipabove=10pt, skipbelow=10pt, backgroundcolor=yellow!8]{warningbox}

% Section formatting
\titleformat{\chapter}[display]
  {\large\bfseries}{\chaptername\ \thechapter}{1em}{\Large\bfseries}
\titleformat{\section}{\normalsize\bfseries}{\thesection.}{0.5em}{}
\titleformat{\subsection}{\normalsize\itshape}{\thesubsection.}{0.5em}{}

% Shorthand macros
\newcommand{\Sadm}{\mathcal{A}}
\newcommand{\Sstate}{\mathcal{S}}
\newcommand{\Inad}{\mathcal{I}}
\newcommand{\Reach}{R(x)}
\newcommand{\OmVol}{\Omega(x)}
\newcommand{\proj}{\pi}
\newcommand{\traj}{\gamma}
\newcommand{\TrajSp}{\Gamma(\mathcal{A})}
\newcommand{\PathAct}{\mathcal{J}[\gamma]}
\newcommand{\PathLag}{L(\gamma,\dot\gamma,t)}
\newcommand{\Efun}{\mathfrak{T}}
\newcommand{\RSVP}{(\Phi,\mathbf{v},S)}
\newcommand{\RepMor}{R_\epsilon}

\pagestyle{fancy}
\fancyhf{}
\fancyhead[LE,RO]{\thepage}
\fancyhead[RE]{\textit{Volume 0: Internal Consistency Document}}
\fancyhead[LO]{\textit{\leftmark}}

\title{%
  \textbf{Volume 0}\\[0.5em]
  \large\textit{Internal Consistency Document}\\[0.3em]
  \normalsize Canonical Symbols, Proof Sketches, Equivalence Briefs,\\
  and the Master Claim Formalization\\[1em]
  \normalsize\textit{Flyxion Fifteen-Volume Curriculum --- Pre-Publication Archive}
}
\author{Flyxion}
\date{June 2026\\[0.3em]
  \small Not for distribution. For internal use only.}

% ================================================================
\begin{document}
% ================================================================

\maketitle
\thispagestyle{empty}

\chapter*{Purpose and Scope}
\addcontentsline{toc}{chapter}{Purpose and Scope}

Volume~0 is not part of the curriculum. It is the document that makes
the curriculum possible to write without introducing contradictions.
It has four components, each corresponding to one class of consistency
problem that, if unresolved, will propagate errors across all fifteen
volumes.

\medskip

\noindent\textbf{Chapter~1: Canonical Symbol Table.} Every object in the
corpus is assigned exactly one symbol. All collisions are resolved.
No volume may introduce a symbol not listed here without amending this
table first.

\noindent\textbf{Chapter~2: Primitive Object Definitions.} Every primitive
object is formally defined. Objects that are derived from primitives
are identified as such. No derived object is treated as primitive
elsewhere in the corpus.

\noindent\textbf{Chapter~3: Asserted Results and Proof Sketches.} Every
result in the corpus that is currently asserted without proof is stated
precisely, assigned an epistemic status, and given a proof sketch,
a proof strategy, or an explicit flag that it is a philosophical proposal
requiring validation before the volume that uses it can be finalized.

\noindent\textbf{Chapter~4: Equivalence Candidate Briefs.} Each of
the nine equivalence candidates identified in the Atlas is given a
one-page formal argument establishing whether it is a genuine
equivalence, a common parent structure, or a shared dynamical form.

\noindent\textbf{Chapter~5: The Master Claim and Path-Action Transition.}
The Master Claim is formalized at two levels (state-functional and
path-variational), the relationship between them is stated precisely,
and the proof obligations for Volume~XV are identified.

\medskip

\begin{warningbox}
\noindent\textbf{Rule}: Writing on any of the fifteen volumes may not
begin until this document has been reviewed and the items flagged as
foundational risks have been addressed. A failed foundational result
changes nearly every volume; a failed technical result changes one
or two. Chapters~3 and 5 identify which is which.
\end{warningbox}

\tableofcontents

% ================================================================
\chapter{Canonical Symbol Table}
% ================================================================

Every symbol in this table is canonical. Where a legacy symbol
conflicts with the canonical assignment, the legacy symbol is
immediately retired for all new writing. Existing documents using
legacy symbols must be updated before submission.

\section{Foundational Geometric Objects}

\medskip
\begin{longtable}{p{3.5cm}p{3cm}p{6.5cm}}
\toprule
\textbf{Concept} & \textbf{Symbol} & \textbf{Notes} \\
\midrule
State space & $\mathcal{S}$ & Full space of all possible configurations.
  Not to be confused with the entropy field $S$. \\
Admissibility space & $\mathcal{A}$ & Sub-manifold of $\mathcal{S}$
  where constraints are satisfied. \\
Admissibility metric & $g$ & Fisher information metric on $\mathcal{A}$.
  In RSVP context, may be extended to $g_\mathrm{RSVP}$. \\
Inadmissibility functional & $\mathcal{I}$ & Non-negative functional;
  $\mathcal{I}(x) = 0 \iff x \in \mathcal{A}$. Not $V$, not $F$,
  not $d^2$ — those are domain-specific instances. \\
Reachability set & $R(x)$ & Set of states reachable from $x$ via
  admissible paths. \\
Reachability volume & $\Omega(x)$ & $\mathrm{Vol}(R(x))$. Measure of
  local future freedom. \\
Constraint family & $C = \{C_i\}$ & $\mathcal{A} = \{x \mid C_i(x)
  \leq 0\;\forall i\}$. \\
\bottomrule
\end{longtable}

\section{Projection and Observation Objects}

\medskip
\begin{longtable}{p{3.5cm}p{3cm}p{6.5cm}}
\toprule
\textbf{Concept} & \textbf{Symbol} & \textbf{Notes} \\
\midrule
Projection (compression) & $\pi: X \to M$ & Smooth surjective map
  from state space to observable manifold. \\
Fiber / preimage & $\pi^{-1}(m)$ & States consistent with observation
  $m$. Epistemic object, not dynamical. \\
Section / reconstruction & $\sigma: M \to X$ & Right inverse of $\pi$;
  $\pi \circ \sigma = \mathrm{id}_M$. Replaces $\mathcal{R}$ for
  reconstruction direction to avoid collision with reachability. \\
General recovery operator & $\mathcal{Q}: M \to X$ & Umbrella for
  all reconstruction operators ($\sigma$, ecphory, Collapse). \\
Observable manifold & $M$ & Target of projection $\pi$. \\
Observation & $m \in M$ & A single observation; $m = \pi(x)$
  for some $x$. \\
Opacity & $\Omega_S + \Omega_M$ & Structural ($\Omega_S$) and
  manufactured ($\Omega_M$) components. Not to be confused with
  reachability volume $\Omega(x)$; subscripts required when both
  appear together. \\
\bottomrule
\end{longtable}

\section{Trajectory and Path Objects}

\medskip
\begin{longtable}{p{3.5cm}p{3cm}p{6.5cm}}
\toprule
\textbf{Concept} & \textbf{Symbol} & \textbf{Notes} \\
\midrule
Admissible trajectory & $\gamma: [0,T] \to \mathcal{A}$ & Primary
  trajectory object; primitive from Volume~IV onward. \\
Trajectory space & $\Gamma(\mathcal{A})$ & Space of all admissible
  trajectories. \\
Path action & $\mathcal{J}[\gamma]$ & $\int_0^T L(\gamma,\dot\gamma,t)
  \, dt$. \\
Path Lagrangian & $L$ & $L: T\mathcal{A} \times [0,T] \to \mathbb{R}$.
  Defined on the tangent bundle of $\mathcal{A}$. \\
Extremal trajectory & $\gamma^*$ & $\operatorname{arg\,ext}_{\gamma
  \in \Gamma(\mathcal{A})} \mathcal{J}[\gamma]$. \\
Repair path & $\gamma_\mathrm{rep}$ & Trajectory satisfying
  $d\Omega/dt > 0$ along its length. \\
Event history & $H_t = (e_1,\ldots,e_n)$ & MEM\textbar{}8 discrete
  trajectory in event space; $H_t \in \Gamma_\mathrm{disc}
  (\mathcal{A})$. \\
\bottomrule
\end{longtable}

\section{RSVP Field Objects}

\medskip
\begin{longtable}{p{3.5cm}p{3cm}p{6.5cm}}
\toprule
\textbf{Concept} & \textbf{Symbol} & \textbf{Notes} \\
\midrule
Scalar field & $\Phi$ & Structural potential / salience / likelihood
  landscape. \\
Vector field & $\mathbf{v}$ & Admissibility flow / information
  transport. \\
Entropy field & $S$ & RSVP dynamical entropy; local accessibility
  density. Distinct from Shannon entropy $H$ and path
  degradation $\sigma_\mathrm{deg}$. \\
RSVP triple & $X = (\Phi, \mathbf{v}, S)$ & Full field state. \\
Lamphron & $\lambda^+$ region & Local maximum of $S$; accessibility
  source. \\
Lamphrodyne & $\lambda^-$ region & Local minimum of $S$; accessibility
  sink; site of structure formation. \\
\bottomrule
\end{longtable}

\section{Operator Objects}

\medskip
\begin{longtable}{p{3.5cm}p{3cm}p{6.5cm}}
\toprule
\textbf{Concept} & \textbf{Symbol} & \textbf{Notes} \\
\midrule
Pop & $\mathsf{Pop}$ & Reduces local entropy at site exceeding
  threshold $\theta_P$. \\
Bind & $\mathsf{Bind}$ & Joins causally compatible events. \\
Collapse & $\mathsf{Collapse}$ & Realization functor $F: \mathcal{H}
  \to \mathcal{C}$; state derived from history. \\
Refuse & $\mathsf{Refuse}$ & Principled exclusion with responsibility
  weight $\rho_x$. \\
Repair morphism & $R_\epsilon$ & Admissibility-preserving map
  reducing $\mathcal{I}$; formal definition pending (see §3.9). \\
Ecphory operator & $\mathcal{E}_\mathrm{cph}$ & Wave-propagation
  retrieval; fires when activation exceeds threshold $\theta_E$. \\
Lamphrodyne operator & $\mathcal{L}_\mathrm{lp}$ & Linear relaxation
  component of RSVP dynamics. \\
Energy functor & $\mathfrak{T}$ & Maps frameworks to energy landscapes;
  see Chapter~5. Not $\mathbb{T}$ (torus). \\
\bottomrule
\end{longtable}

\section{Resolved Symbol Conflicts}

The following conflicts existed in the prior corpus. The canonical
resolution is mandatory in all new and revised writing.

\medskip
\begin{longtable}{p{3cm}p{4cm}p{4cm}p{3cm}}
\toprule
\textbf{Concept} & \textbf{Retired symbols} & \textbf{Canonical} &
\textbf{Reason} \\
\midrule
Entropy (three kinds) & $S$ (all three uses) & $S$ for RSVP field;
  $H$ for Shannon; $\sigma_\mathrm{deg}$ for path degradation &
  Three genuinely different objects \\
Inadmissibility & $V$, $F$, $d^2$, $\mathcal{O}$ & $\mathcal{I}$
  universally & Domain-specific names are instances \\
Reachability (set vs.\ volume) & $\Omega$, $\mathcal{F}_\mathrm{soc}$
  for same concept & $R(x)$ for set; $\Omega(x)$ for volume;
  $\mathcal{F}_\mathrm{soc}$ for social freedom & Prevents
  conflation \\
Reconstruction direction & $\mathcal{R}$, $F$, $\sigma$ & $\sigma$
  for section; $\mathcal{Q}$ for general recovery & $\mathcal{R}$
  collides with reachability \\
Free energy (Bayesian) & $F$, $\mathcal{I}$ & $F$ reserved for
  variational free energy in Bayesian contexts; $\mathcal{I}$ for
  inadmissibility elsewhere & Prevents thermodynamic confusion \\
Energy functor & $\mathbb{T}$, $\mathscr{T}$ & $\mathfrak{T}$ &
  $\mathbb{T}$ is the torus; $\mathscr{T}$ is non-standard \\
\bottomrule
\end{longtable}

% ================================================================
\chapter{Primitive Object Definitions}
% ================================================================

This chapter defines every primitive object in the corpus. Derived
objects are defined in terms of primitives and labeled as such.
No derived object is to be treated as foundational elsewhere.

\section{The Admissibility Structure}

\begin{definition}[Admissibility Structure]
An \emph{admissibility structure} is a quadruple
$(\mathcal{S}, \mathcal{A}, g, \mathcal{I})$ where:
\begin{itemize}
\item $\mathcal{S}$ is a smooth manifold (the state space).
\item $\mathcal{A} \subseteq \mathcal{S}$ is a sub-manifold defined
  by $\mathcal{A} = \{x \in \mathcal{S} \mid C_i(x) \leq 0\;\forall i
  \in I\}$ for a constraint family $\{C_i\}_{i \in I}$.
\item $g$ is a Riemannian metric on $\mathcal{A}$; the canonical
  choice is the Fisher information metric
  $g_{jk} = \mathbb{E}[\partial_j \log p\, \partial_k \log p]$.
\item $\mathcal{I}: \mathcal{S} \to \mathbb{R}_{\geq 0}$ is a
  smooth inadmissibility functional with $\mathcal{I}(x) = 0
  \iff x \in \mathcal{A}$ and $\mathcal{I}(x) > 0$ otherwise.
\end{itemize}
\end{definition}

\begin{remark}
The admissibility structure is the common parent of several objects
that appear under different names across frameworks. These are
\emph{not} equivalent to each other; they are \emph{derived from}
the admissibility structure by additional specifications:
$R(x) \subseteq \mathcal{A}$ (dynamical subobject),
$\pi^{-1}(m) \subseteq X$ (epistemic subobject),
$\mathcal{P}$ in MEM\textbar{}8 (historical subobject). See
Chapter~4, Candidate~1.
\end{remark}

\section{The Projection Structure}

\begin{definition}[Projection Structure]
A \emph{projection structure} is a triple $(X, M, \pi)$ where
$X$ is a state space, $M$ is an observable manifold, and
$\pi: X \to M$ is a smooth surjective map (a projection or
submersion). The \emph{fiber} over $m \in M$ is
$\pi^{-1}(m) = \{x \in X \mid \pi(x) = m\}$.
\end{definition}

\begin{definition}[Operational Fiber Death]
A state $m \in M$ exhibits \emph{operational fiber death} when
$\pi^{-1}(m) \neq \varnothing$ but the measure of admissible
reconstruction pathways satisfies $\mu(\mathcal{P}_m) \to 0$.
That is: the preimage is non-empty in principle but unreachable
in practice.
\end{definition}

\section{The Admissible Trajectory}

\begin{definition}[Admissible Trajectory]
\label{def:traj}
An \emph{admissible trajectory} is a curve
$\gamma: [0,T] \to \mathcal{A}$ such that:
\begin{enumerate}
\item $\gamma(t) \in \mathcal{A}$ for all $t \in [0,T]$.
\item $\dot\gamma(t) \in T_{\gamma(t)}\mathcal{A}$ wherever
  $\partial\mathcal{A}$ is non-empty.
\end{enumerate}
The \emph{trajectory space} is $\Gamma(\mathcal{A}) = \{\gamma:
[0,T] \to \mathcal{A} \mid \gamma\text{ admissible}\}$ with the
$C^1$ topology.

A \emph{discrete admissible trajectory} is a finite sequence
$(\gamma_0, \gamma_1, \ldots, \gamma_n) \in \mathcal{A}^{n+1}$
where consecutive pairs $(\gamma_k, \gamma_{k+1})$ are connected
by a short admissible path.
\end{definition}

\begin{remark}
The following objects across the corpus are all specializations
of Definition~\ref{def:traj}: MEM\textbar{}8 event history $H_t$
(discrete trajectory in event space), repair path $\gamma_\mathrm{rep}$
(trajectory reducing $\mathcal{I}$ along its length), agency plan
(trajectory maximizing $\mathbb{E}[\Omega(\gamma(T))]$),
ecphoric path (trajectory from present state to target memory),
Yarncrawler route (trajectory completing partial world-state),
TARTAN trajectory (trajectory preserved under decomposition).
\end{remark}

\section{The Path Lagrangian}

\begin{definition}[Path Lagrangian]
A \emph{path Lagrangian} is a smooth function
$L: T\mathcal{A} \times [0,T] \to \mathbb{R}$
on the tangent bundle of $\mathcal{A}$. The associated
\emph{path action} is
\[
\mathcal{J}[\gamma] = \int_0^T L(\gamma(t), \dot\gamma(t), t)\, dt.
\]
An \emph{extremal trajectory} $\gamma^*$ satisfies the
Euler-Lagrange equations:
\[
\frac{d}{dt}\frac{\partial L}{\partial \dot\gamma^i}
- \frac{\partial L}{\partial \gamma^i} = 0.
\]
\end{definition}

\begin{definition}[Canonical Lagrangian Instantiations]
The following are the specific Lagrangians used by path-action
frameworks in the corpus:
\begin{align*}
L_\mathrm{mech}(\gamma,\dot\gamma) &=
  \tfrac{1}{2}g_{ij}\dot\gamma^i\dot\gamma^j - \mathcal{I}(\gamma)
  &\text{(recovers gradient flow)}\\
L_\mathrm{visc}(\gamma,\dot\gamma) &=
  \mathcal{I}(\gamma) + \lambda\|\dot\gamma\|^2_\eta
  &\text{(MEM\textbar{}8, ecphory with viscosity)}\\
L_\mathrm{rep}(\gamma,\dot\gamma) &=
  d(f(\gamma), f^*)^2 - \eta\Omega(\gamma)
  &\text{(repair Lagrangian, Volume IX)}\\
L_\mathrm{SP}(\gamma, e_t) &=
  -\log p(e_t \mid H_t)
  &\text{(Spherepop event log-likelihood)}
\end{align*}
\end{definition}

\section{The Repair Morphism (Provisional Definition)}

\begin{definition}[Repair Morphism --- Provisional]
\label{def:repair}
A map $R_\epsilon: \mathcal{S} \to \mathcal{S}$ is a
\emph{repair morphism} if it satisfies:
\begin{enumerate}
\item \textbf{Admissibility}: $R_\epsilon(x) \in \mathcal{A}$ for
  all $x$ in a neighborhood of $\mathcal{A}$.
\item \textbf{Inadmissibility reduction}:
  $\mathcal{I}(R_\epsilon(x)) \leq \mathcal{I}(x)$ with equality
  only if $x \in \mathcal{A}$.
\item \textbf{Path consistency}: There exists an admissible
  trajectory $\gamma: [0,1] \to \mathcal{A}$ with $\gamma(0) = x$
  and $\gamma(1) = R_\epsilon(x)$.
\end{enumerate}
\end{definition}

\begin{warningbox}
\textbf{Foundational gap}: Definition~\ref{def:repair} is provisional.
The repair morphism is characterized by $d\Omega/dt > 0$ along its
trajectory, but this characterization is a candidate proposition
(§3.9), not a theorem. The definition above gives conditions; a
satisfying construction that meets them requires either: (a) proof
that $L_\mathrm{rep}$ above has well-defined extremals in
$\Gamma(\mathcal{A})$, or (b) an explicit construction of
$R_\epsilon$ as a flow map. Volume~VIII (Repair Theory) cannot be
finalized until one of these is provided.
\end{warningbox}

% ================================================================
\chapter{Asserted Results and Proof Sketches}
% ================================================================

\noindent\textbf{Epistemic status codes:}
\textbf{Pr} = proved in existing document;
\textbf{Sk} = proof sketch exists;
\textbf{Co} = conjecture with identified proof strategy;
\textbf{Nu} = numerical evidence only;
\textbf{As} = asserted without proof strategy;
\textbf{Ph} = philosophical proposal requiring validation;
\textbf{Ga} = gap, not yet precisely stated.

\bigskip
\noindent Results marked \textbf{As} or \textbf{Ph} that appear in
the \emph{foundational gap} category must be addressed before writing
the volumes that depend on them. Results marked \textbf{Ga} must be
precisely stated before appearing in any volume.

\section{Volume I: Constraint Geometry}

\begin{theorem}[Projection-Collapse Principle, Weak Form] \textbf{(Pr)}
Let $\pi: X \to M$ be a smooth projection and let
$f: \pi^{-1}(U) \to \mathbb{R}$ be a function on a fiber over
an open $U \subseteq M$. Then information loss under $\pi$
satisfies $H(X) - H(\pi(X)) \geq 0$ with equality iff $\pi$
is bijective on $U$.
\end{theorem}
\begin{proof}[Proof sketch]
Data-processing inequality applied to the Markov chain
$X \to \pi(X)$. Equality condition follows from injectivity.
\end{proof}

\begin{conjecture}[Projection-Collapse Principle, Strong Form] \textbf{(Co)}
Let $\pi: X \to M$ have positive-dimensional fibers. Then
the reconstruction operator $\mathcal{Q}: M \to X$ satisfying
$\pi \circ \mathcal{Q} = \mathrm{id}_M$ is generically ill-posed:
small perturbations in $m$ produce large perturbations in
$\mathcal{Q}(m)$.
\end{conjecture}
\begin{proof}[Proof strategy]
Apply Hadamard's three conditions for well-posedness. Non-injectivity
of $\pi$ implies non-uniqueness of $\mathcal{Q}$, violating the
second condition. Continuity failure follows from fiber
non-compactness or high curvature. Tikhonov regularization provides
a stable approximation, formalizing $\mathcal{Q}_\lambda$ as a
regularized inverse.
\end{proof}

\begin{proposition}[Constraint-Information Equivalence] \textbf{(Pr)}
Let $\mathcal{S}$ be a finite state space of size $N$ and
$\mathcal{A} \subseteq \mathcal{S}$ the admissible subset of size
$k$. The information encoded by the constraint family
$\{C_i\}$ is $\log N - \log k = \log(N/k)$.
\end{proposition}

\begin{theorem}[Bishop-Gromov Comparison for Admissibility] \textbf{(Pr)}
Let $(\mathcal{A}, g)$ have Ricci curvature bounded below by
$(n-1)\kappa$. Then the volume of metric balls
$B(x,r) \subset \mathcal{A}$ satisfies
$\mathrm{Vol}(B(x,r)) \leq V_\kappa(r)$
where $V_\kappa(r)$ is the volume of a ball of radius $r$ in the
model space of constant curvature $\kappa$.
\end{theorem}
\begin{proof}[Proof sketch]
Standard Bishop-Gromov, applied to the Fisher information metric
on $\mathcal{A}$. Requires $\mathcal{A}$ to be complete as a
Riemannian manifold.
\end{proof}

\section{Volume II: RSVP Field Theory}

\begin{claim}[RSVP Second Law] \textbf{(As)}
For solutions $(\Phi, \mathbf{v}, S)$ of the RSVP field equations,
global accessibility satisfies
$\frac{d}{dt}\int_\Omega S\,dx \leq 0$
under appropriate boundary conditions.
\end{claim}
\begin{proof}[Proof obligation]
Integrate the $S$-equation $\partial_t S = D_S\nabla^2 S +
R(\Phi,\mathbf{v},S)$ over $\Omega$. The Laplacian term vanishes
with Neumann boundary conditions. A sufficient condition is
$\int_\Omega R(\Phi,\mathbf{v},S)\,dx \leq 0$, which constrains the
reaction term. The physical interpretation is that lamphrodyne
processes are net entropy-consuming. \textbf{Status: must be proved
before Volume~XIII can be finalized.}
\end{proof}

\begin{claim}[Lamphrodyne Attractor Stability] \textbf{(As)}
A lamphrodyne region (local minimum of $S$) is Lyapunov stable
under the RSVP dynamics.
\end{claim}
\begin{proof}[Proof obligation]
Construct a Lyapunov functional $V[\Phi,\mathbf{v},S]$ that decreases
along RSVP trajectories near a lamphrodyne configuration. The
candidate is $V = \int_\Omega |\nabla S|^2 + \alpha(\Phi - \Phi_0)^2
+ \beta|\mathbf{v}|^2\,dx$. Show $\dot V \leq 0$. \textbf{Status:
must be proved before Volumes~XIII and~XIV.}
\end{proof}

\begin{conjecture}[RSVP-RG Equivalence] \textbf{(Co)}
The coarse-graining flow induced by lamphrodyne relaxation is
equivalent, in appropriate limits, to Wilsonian renormalization
group flow on the space of effective field theories.
\end{conjecture}
\begin{proof}[Proof strategy]
Show that integrating out high-frequency modes in the RSVP path
integral $Z = \int \mathcal{D}\Phi\,\mathcal{D}\mathbf{v}\,\mathcal{D}S
\, e^{-\mathcal{A}[X]}$ produces a flow equation on the coupling
constants of the effective action equivalent to the Wetterinck
equation (functional RG). The entropy diffusion coefficient $D_S$
plays the role of the RG scale parameter.
\end{proof}

\section{Volume III: Projection and Inference}

\begin{definition}[CLIO Intelligence]
A system is \emph{CLIO-intelligent} to degree $\alpha$ at
observation $m$ if it can identify a subset
$\mathcal{A}(m) \subseteq \pi^{-1}(m)$ satisfying
$|\mathcal{A}(m)| \leq \alpha |\pi^{-1}(m)|$
using only constraints derivable from context and memory.
\end{definition}

\begin{claim}[CLIO Inverse Problem Ill-Posedness] \textbf{(Ga)}
The reconstruction problem of recovering $x \in \pi^{-1}(m)$ from
observation $m$ alone is generically ill-posed in the sense of Hadamard.
\end{claim}
\begin{proof}[Proof obligation]
State precisely: for which class of projections $\pi$ and which
topology on $X$ is the problem ill-posed? Then prove the three
conditions fail. Note: the Projection-Collapse strong form (§3.1)
is the precursor result; this is a specialization to the inference
setting. \textbf{Gap: precise statement needed before Volume~III
is drafted.}
\end{proof}

\begin{theorem}[Reconstruction Stability] \textbf{(Pr)}
The Tikhonov-regularized reconstruction $\mathcal{Q}_\lambda(m) =
\operatorname{arg\,min}_{x \in X}\{\|x\|^2 + \lambda^{-1}
\|\pi(x) - m\|^2\}$ satisfies an $L^2$ stability estimate:
$\|\mathcal{Q}_\lambda(m_1) - \mathcal{Q}_\lambda(m_2)\|
\leq C_\lambda\|m_1 - m_2\|$ for a constant $C_\lambda$ depending on
$\lambda$ and the spectral properties of $\pi^*\pi$.
\end{theorem}

\begin{theorem}[Operational Fiber Death] \textbf{(Pr)}
Let $\mathcal{P}_m$ denote the set of admissible reconstruction
pathways to the fiber $\pi^{-1}(m)$. If $\mu(\mathcal{P}_m) \to 0$
while $\pi^{-1}(m) \neq \varnothing$, the state $m$ exhibits
operational fiber death: it is theoretically reachable but
practically unrecoverable.
\end{theorem}

\begin{claim}[Admissible Stabilization --- AI Convergence Certification]
\textbf{(As)}
\label{claim:ai-cert}
Let $s_t \in M$ be the observable state of a computational system
at time $t$, let $x_t \in X$ be the latent state with
$\pi(x_t) = s_t$, and let $\mathcal{A} \subseteq X$ be the
admissibility space encoding safety constraints. Define:
\begin{itemize}
\item $d_M(s_t, s_{t-1})$: metric distance between consecutive
  observable states (motion in operational space);
\item $C_t \in [0,1]$: corrected consistency score measuring
  agreement of $s_t$ against an independent prior or prospective
  candidate (not self-referential);
\item $E_\pi(s_t)$: surrogate error, measuring agreement across
  independent projections of the latent state.
\end{itemize}
Then the joint condition:
\[
d_M(s_t, s_{t-1}) \to 0
\;\;\text{AND}\;\;
C_t \to 1
\;\;\text{AND}\;\;
E_\pi(s_t) \leq \eta
\;\;\Longrightarrow\;\;
P\!\left(\mathrm{Stable}_{\mathcal{A}}(x_t)\right) \to 1
\]
constitutes \emph{admissible stabilization}: the system has
converged to an admissible latent state with high probability.
\end{claim}

\begin{proof}[Proof obligation and theoretical context]
The claim addresses the \emph{Ouroboros trap}: a system that
measures only $d_M(s_t, s_{t-1}) \to 0$ certifies convergence
in observable space $M$ without confirming convergence in latent
space $X$. Since $\pi$ is non-injective, a system can refine an
incorrect plan into a perfectly self-consistent incorrect plan ---
converging in the projection without converging in the fiber.

The three-condition conjunction is designed to break this trap:
\begin{enumerate}
\item $d_M \to 0$ (Ouroboros signal): motion ceases in operational
  space. Necessary but not sufficient.
\item $C_t \to 1$ (CLIO condition): real comparison against an
  independent candidate, not a tautology. Catches shadow-space
  convergence.
\item $E_\pi \leq \eta$ (surjective consistency): agreement across
  multiple independent projections $\{\pi_k\}$ of the latent state
  confirms that the convergence is not an artifact of one projection.
\end{enumerate}

\textbf{What must be proved}: (a) that the conjunction is sufficient
for $P(\mathrm{Stable}_\mathcal{A}(x_t)) \to 1$ under appropriate
conditions on $\pi$ and the dynamics; (b) that $C_t$ is computable
without access to $x_t$ directly; (c) that the three conditions are
independent (no two are jointly sufficient without the third).

\textbf{Source}: Epistemic Architecture slides, Slide~5.
\textbf{Status}: \textbf{As}. The three conditions are identified
and motivated; no formal proof of sufficiency exists.

\textbf{Why it matters}: This result, if proved, provides the first
formally grounded AI safety certification criterion derived from
CLIO projection theory. It appears in Volume~III as a worked
application of the CLIO pipeline. It also connects to the Invariant
Discovery Principle: what the three-condition test certifies is
that the structural invariant $\mathcal{I}(T_i(C)) = \mathcal{I}(C)$
is preserved across observable transformations, not merely that
the surface has stopped moving.
\end{proof}

\begin{remark}[The Projection Trap and Sealed Systems]
The Ouroboros trap is the computational instance of a more general
failure mode identified across the corpus: a sealed interpretive
system that progressively eliminates reachable futures by
reinterpreting all evidence as confirmation of its current state.
In the Szukalski analysis (Epistemic Blueprint), Zermatism
exemplifies this: every disconfirming piece of evidence is
reinterpreted as evidence of the debasement the system diagnoses.
In AI systems, the analogous failure is a trajectory classifier
that evaluates $d(s_t, s_{t-1}) \to 0$ while remaining blind to
the latent space structure it is converging on. The three-condition
certification is the formal antidote to sealed-system convergence
in computational contexts.
\end{remark}

\begin{claim}[MEM|8 Reconstruction Theorem] \textbf{(Ga)}
Let $H_n = (e_1,\ldots,e_n)$ be a MEM\textbar{}8 event history
and $s_n = \bigcap_i \mathcal{C}(e_i)$ the reconstructed state.
Under what conditions on the event algebra is $s_n$ uniquely
determined by $H_n$, and when is the reconstruction continuous
with respect to perturbations of the event sequence?
\end{claim}
\begin{proof}[Proof obligation]
This is the MEM\textbar{}8 analog of the Reconstruction Stability
theorem. The reconstruction is intersection of constraint sets; its
well-posedness depends on the intersection being non-empty and
having non-degenerate geometry. \textbf{Gap: precise statement and
proof needed before Volume~V is finalized.}
\end{proof}

\begin{claim}[Ecphoric Threshold Existence] \textbf{(As)}
For any memory state $m^* \in \Gamma(\mathcal{A})$ and any present
state $x_0 \in \mathcal{A}$, there exists a threshold $\theta_E > 0$
such that the ecphoric retrieval wave initiated at $x_0$ reaches
$m^*$ if and only if the wave amplitude at $x_0$ exceeds $\theta_E$.
\end{claim}
\begin{proof}[Proof obligation]
Model ecphory as a damped wave equation on $\mathcal{A}$ with
initial condition at $x_0$. The threshold $\theta_E$ is determined
by the viscosity $\eta(x,t)$ of the semantic manifold between $x_0$
and $m^*$. Existence requires the wave equation to have a unique
solution; the threshold exists if the solution operator is continuous.
\textbf{Must be proved before Volumes~V and~VIII.}
\end{proof}

\begin{claim}[HYDRA Cognitive Emergence] \textbf{(As)}
In a HYDRA system with modules $\{f_i\}_{i=1}^N$ and coordination
penalty $\Omega(\mathbf{z})$, complex cognitive behaviors emerge
from coordination without any individual module possessing those
behaviors.
\end{claim}
\begin{proof}[Proof obligation]
Formalize "emergence" as: behavior $B$ is emergent if it is
expressed by the system's aggregate output but not by any individual
module's output restricted to its input domain. Show that the
coordination penalty $\Omega$ induces global coherence not implied
by any local $f_i$. This requires a formal stability analysis of the
coupled system $\dot{\mathbf{z}} = \sum_a W_a f_a(\mathbf{z}_a)
- \nabla_\mathbf{z}\Omega(\mathbf{z})$. \textbf{Gap: stability
analysis is the central missing piece in Volume~VII.}
\end{proof}

\begin{claim}[HYDRA Stability] \textbf{(Ga)}
The coupled HYDRA system with $N$ modules converges to a coherent
equilibrium $\mathbf{z}^*$ satisfying $\sum_a W_a f_a(\mathbf{z}^*_a)
= \nabla_\mathbf{z}\Omega(\mathbf{z}^*)$ under conditions on the
coupling weights $\{W_a\}$ and the coordination penalty $\Omega$.
\end{claim}
\begin{proof}[Proof obligation]
Construct a Lyapunov function for the coupled system. Natural
candidate: $V(\mathbf{z}) = \Omega(\mathbf{z}) +
\sum_a\frac{W_a}{2}\|f_a(\mathbf{z}_a)\|^2$. Show $\dot V \leq 0$.
Identify sufficient conditions on $\Omega$ (e.g., strong convexity)
for global convergence. \textbf{Must be proved before Volume~VII.}
\end{proof}

\begin{claim}[Repair Characterization Candidate] \textbf{(As)}
The condition $d\Omega(\gamma(t))/dt > 0$ along an admissible
trajectory $\gamma$ is equivalent to the trajectory satisfying the
repair Lagrangian Euler-Lagrange equations for
$L_\mathrm{rep} = d(f(\gamma),f^*)^2 - \eta\Omega(\gamma)$.
\end{claim}
\begin{proof}[Proof obligation]
Compute the Euler-Lagrange equations for $L_\mathrm{rep}$ and show
they imply $d\Omega/dt > 0$ along solutions. This requires
(a) $\Omega$ to be differentiable along $\gamma$, and (b) the
coupling $\eta > 0$ to be sufficiently large. \textbf{This is the
central proof obligation for Volume~VIII. Resolve before drafting.}
\end{proof}

\section{Volume VII: Spherepop}

\begin{claim}[Spherepop Confluence] \textbf{(Ga)}
The Spherepop reduction system is confluent: for any initial
configuration $C_0$, all maximal reduction sequences terminate
at the same normal form up to trace-equivalence.
\end{claim}
\begin{proof}[Proof obligation]
Apply Newman's lemma: confluence follows from local confluence and
strong normalization. (1) Strong normalization: show that every
Spherepop reduction sequence terminates (no infinite chains). Natural
measure: the total entropy $\int S\,dx$ of the associated RSVP
state is non-increasing under Pop, and bounded below by zero.
(2) Local confluence: show that if $C_0 \to C_1$ and $C_0 \to C_2$
by single steps, then $C_1$ and $C_2$ have a common reduct.
Check this for all pairs of rules. \textbf{Must be proved before
Volume~IX is finalized.}
\end{proof}

\section{Volume X: Coordination Geometry}

\begin{claim}[Kuramoto Reduction from RSVP] \textbf{(Co)}
In the limit where $\Phi$ and $S$ are slowly varying compared to
$\mathbf{v}$, the RSVP vector field equation reduces to a
generalized Kuramoto model $\dot\theta_i = \omega_i +
\sum_j K_{ij}\sin(\theta_j - \theta_i)$ where $\theta_i$ are
phases of the $\mathbf{v}$ field and $K_{ij}$ are coupling
constants derived from $\nabla S$.
\end{claim}
\begin{proof}[Proof strategy]
Project the RSVP $\mathbf{v}$-equation onto a phase representation
$v_i = A_i e^{i\theta_i}$ using Stuart-Landau normal form near
a Hopf bifurcation. The entropy gradient $\gamma\nabla S$ becomes
the coupling term. Amplitude variations decouple to leading order,
leaving a pure phase equation. The coupling strength $K_{ij}$ depends
on the overlap integrals of the entropy gradient with the phase
mode functions. Requires the RSVP $\mathbf{v}$-equation to have
a Hopf bifurcation at the relevant parameter values.
\end{proof}

\section{Foundational Results: High Priority}

These results are foundational in the sense that their failure
restructures nearly all volumes. They are listed separately from
the technical gaps above.

\begin{claim}[Reachability Principle] \textbf{(Sk)}
\label{claim:rp}
The effective freedom of an agent or system at state $x$ is
proportional to the reachability volume $\Omega(x) =
\mathrm{Vol}(R(x))$. Formally, the optimal policy for any
persistence-oriented objective is equivalent to maximizing
$\Omega$ along the system's trajectory.
\end{claim}
\begin{proof}[Proof sketch --- from Institutional Attractors, Appendix A]
Define the local contraction rate
$\Lambda(x,t) = -\frac{d}{dt}\log\mathcal{V}(x,t)$
where $\mathcal{V}(x,t) = \int_{R(x,t)}d\mu$ is the reachable
future volume. If $\Lambda(x,t) > 0$ over $[t_0,t_1]$, then:
\[
\mathcal{V}(t_1) = \mathcal{V}(t_0)\exp\!\left(
  -\int_{t_0}^{t_1}\Lambda(\tau)\,d\tau\right).
\]
Persistent reachability closure therefore produces exponential
loss of admissible futures. A persistence-oriented policy must
hold $\Lambda(x,t) \leq 0$ along its trajectory, which is
equivalent to non-decreasing $\Omega$. The Bellman formulation
follows by setting the value function $V(x) = \log\Omega(x)$.

\textbf{Status upgraded from As to Sk.} The contraction-rate
argument is in Appendix~A of the Institutional Attractors essay.
What remains unproved: (a) that $\Lambda \leq 0$ is sufficient
as well as necessary for persistence, and (b) that maximizing
$\Omega$ is optimal across all persistence-objective classes, not
only those that reduce to reachability volume.
\textbf{If false, affects Volumes IV, VI, VIII, XII, XV.}
\end{proof}

\begin{claim}[Reachability-Inadmissibility Candidate Identity] \textbf{(As)}
\label{claim:ident}
$\mathcal{I}(x) = -\log\Omega(x)$.
\end{claim}
\begin{proof}[Proof obligation]
This is a candidate definition, not a theorem. Two steps required:
(1) Show that $-\log\Omega(x)$ satisfies the defining properties
of $\mathcal{I}$: non-negative, zero iff $x \in \mathcal{A}$, smooth.
For $\Omega$ this requires $R(x) = \mathcal{A}$ when $x \in \mathcal{A}$
(which is true) and $R(x) \subsetneq \mathcal{A}$ otherwise.
(2) Verify consistency with the RSVP entropy field $S$. Since $S$
encodes local accessibility density and $\Omega(x)$ is global
reachability volume, they are related but not identical. The
precise relationship is: $S(x) \sim -\nabla_x\log\Omega(x)$
(accessibility gradient equals entropy gradient) is a conjecture
requiring verification. \textbf{Until verified, present as
candidate identity, not theorem.}
\end{proof}

\begin{claim}[Optimal Opacity] \textbf{(Sk)}
The fitness functional $F(\Omega_M) = F_0 + \kappa\Omega_M -
\eta\Omega_M^2$ has a maximum at $\Omega_M^* = \kappa/2\eta$,
predicting that institutions under opacity selection converge to
locally optimal rather than maximal opacity.
\end{claim}
\begin{proof}[Proof sketch --- from Institutional Attractors, §7.1]
The optimum follows directly: $F'(\Omega_M) = \kappa - 2\eta\Omega_M
= 0$ gives $\Omega_M^* = \kappa/2\eta$, and $F''(\Omega_M^*)
= -2\eta < 0$ confirms it is a maximum. The systemic fragility
corollary ($\Omega_M > \Omega_M^*$ implies internal coordination
costs dominate) follows from the concavity. The 2008 financial
crisis is cited as an instance: securitization opacity exceeded
the institutional coordination threshold, making risk-modeling
impossible from within.

\textbf{Status upgraded from As to Sk.} The optimum derivation
is complete. What remains unproved is the derivation of the
quadratic fitness functional from first principles:
(a) the linear term $\kappa\Omega_M$ must be derived from the
replicator dynamics as the fitness advantage of external
impermeability; (b) the quadratic term $-\eta\Omega_M^2$ must
be derived from coordination-cost theory as the penalty for
internal opacity. Without this derivation, the functional form
is an ansatz, not a theorem. Provide derivation before
Volume~IV (Simulated Agency / Institutional Attractors
application) is finalized.
\end{proof}

% ================================================================
\chapter{Equivalence Candidate Briefs}
% ================================================================

Each candidate receives a one-page formal brief establishing
its status: \textbf{GE} (genuine equivalence), \textbf{CP}
(common parent), or \textbf{SD} (shared dynamical form).

\section{Candidate 1: The Admissibility Cluster}

\begin{flybox}
\textbf{Objects}: $\mathcal{A}$, $R(x)$, $\pi^{-1}(m)$,
$\mathcal{P}$ (MEM\textbar{}8), $\Omega_t$ (Spherepop).

\textbf{Verdict}: \textbf{CP} --- Common Parent Structure.

\textbf{Argument}: The objects occupy different categorical levels.
$R(x) \subseteq \mathcal{A}$: reachability is a dynamical subobject
depending on a starting point. $\pi^{-1}(m) \subseteq X$: a fiber
is an epistemic subobject depending on an observation. $\mathcal{P}$
in MEM\textbar{}8: a historical subobject depending on an event
sequence. $\Omega_t$ in Spherepop: a probability distribution, not
a set. These are not the same object, and defining them identically
would confuse dynamical accessibility, epistemic ambiguity, and
historical possibility.

\textbf{Correct claim}: All are generated from the admissibility
structure $(\mathcal{S}, \mathcal{A}, g, \mathcal{I})$ by additional
specifications. Volume~I defines the tuple; subsequent volumes specify
which projection $\pi$, which starting state, and which event history
they are working with.

\textbf{Consequence}: No merger. Volume~I defines the parent structure;
each downstream volume defines its specialized subobject by
explicit construction.
\end{flybox}

\section{Candidate 2: The Projection Cluster}

\begin{flybox}
\textbf{Objects}: $\pi: X \to M$, rendering operator $\mathcal{R}$
(CLIO), realization functor $F: \mathcal{H} \to \mathcal{C}$
(Spherepop), coarse-graining map (RSVP-RG), compression maps
(Semantic Infrastructure).

\textbf{Verdict}: \textbf{CP} --- Common Parent, with important
directionality distinction.

\textbf{Argument}: The compression-direction maps (CLIO rendering,
RSVP coarse-graining, semantic compression) are all instances of
$\pi: X \to M$. The Spherepop realization functor goes in the
opposite direction: from history (small) to state (large). It is
closer to a \emph{section} $\sigma: M \to X$ than to $\pi$.

\textbf{Correct claim}: All compression maps are instances of $\pi$.
All reconstruction/realization maps are instances of $\sigma$ or the
general recovery operator $\mathcal{Q}$. The two directions must
not be conflated.

\textbf{Consequence}: No merger. Use $\pi$ for compression, $\sigma$
for section, $\mathcal{Q}$ for general recovery.
\end{flybox}

\section{Candidate 3: The Tension / Free Energy Cluster}

\begin{flybox}
\textbf{Objects}: $V(x,t)$ (Institutional Attractors), $F(q,o)$
(Simulated Agency), $d(f(x),f^*)^2$ (Repair), $\Omega(\mathbf{z})$
(HYDRA incoherence), $\mathcal{O}_\mathrm{merge}$ (Semantic
Infrastructure).

\textbf{Verdict}: \textbf{SD} --- Shared Dynamical Form, pending
verification of Energy Functor.

\textbf{Argument}: Every object in this cluster is the potential
$\mathcal{I}$ in the gradient flow master equation for its respective
framework. The formal spine document confirms this: every volume's
master equation has the form $\dot{x} = -\nabla\mathcal{I} +
\text{terms}$. This is more than analogy: the gradient flow structure
is genuinely shared.

\textbf{Pending}: Whether the individual $\mathcal{I}$ objects are
related by coordinate change (genuine equivalence) or are genuinely
independent potentials on genuinely different state spaces (shared
form only) depends on whether the Energy Functor $\mathfrak{T}$
exists and is faithful (see Candidate~9). If $\mathfrak{T}$ is
confirmed, this cluster upgrades to \textbf{GE}.

\textbf{Consequence}: Use $\mathcal{I}$ as the canonical symbol
universally. Domain-specific names ($V$, $F$, $d^2$) become
instantiation labels, not independent primitives.
\end{flybox}

\section{Candidate 4: The Collapse / Reconstruction Cluster}

\begin{flybox}
\textbf{Objects}: Reconstruction $\mathcal{Q}_\lambda$ (Vol.~I),
ecphory $\mathcal{E}_\mathrm{cph}$ (cognitive), Collapse
$\mathsf{Collapse}$ (Spherepop), MEM\textbar{}8 state recovery
$s_n = \bigcap_i \mathcal{C}(e_i)$, repair morphism $R_\epsilon$.

\textbf{Verdict}: \textbf{CP} --- Common Parent with different
optimality criteria.

\textbf{Argument}: All are instances of the general recovery
operator $\mathcal{Q}: M \to X$ defined as a right inverse of
$\pi$ under some optimality criterion. The criterion determines
the specific operator:
\begin{itemize}
\item Tikhonov: minimize $\|x\|^2$ subject to $\pi(x) \approx m$
\item Ecphory: maximize activation given viscosity-damped wave propagation
\item Spherepop Collapse: functional composition
  $T_{x_n} \circ \cdots \circ T_{x_1}(s_0)$
\item MEM\textbar{}8: constraint intersection
  $\bigcap_i \mathcal{C}(e_i)$
\end{itemize}
Note: the repair morphism $R_\epsilon$ is \emph{not} an instance of
$\mathcal{Q}$. It operates on broken states in $\mathcal{S}$, not
on observations in $M$. Repair does not invert a projection; it
corrects a trajectory. Keep separate.

\textbf{Consequence}: Define $\mathcal{Q}$ as a primitive with
criterion as parameter. Each framework specifies its criterion.
\end{flybox}

\section{Candidate 5: The Reachability Cluster}

\begin{flybox}
\textbf{Objects}: Reachability volume $\Omega(x)$, social freedom
$\mathcal{F}_\mathrm{soc} = \sum_i\mathrm{Vol}(\mathcal{A}_i)$,
simulated reachability $\mathcal{R}(q)$, repair bound
$\eta\mathcal{R}(x)$, diagnostic volume $V(t)$ (Self-Concealment),
geometric invariance $\mathcal{G}(Q)$ (consciousness).

\textbf{Verdict}: \textbf{SD} --- Shared Dynamical Form.

\textbf{Argument}: All are instances of $\Omega(x) =
\mathrm{Vol}(R(x))$ in different state spaces. Social freedom
measures reachability in social configuration space. The diagnostic
volume $V(t)$ measures reachability in the space of epistemological
positions (the vantage points from which the mechanism can be
identified). Geometric invariance $\mathcal{G}(Q)$ is conjectured
to measure reachability in quale-state space.

\textbf{Consequence}: Use $\Omega$ with appropriate subscripts
identifying the state space. Unify under the reachability field
$\mathfrak{R}(x)$ as a primitive if Candidate~8 is confirmed.
\end{flybox}

\section{Candidate 6: Cognitive Branch Sequence}

\begin{flybox}
\textbf{Objects}: CLIO, MEM\textbar{}8, Simulated Agency, HYDRA.

\textbf{Verdict}: \textbf{CP} --- Common Parent, organized as
coherent four-volume sequence, not merged into one volume.

\textbf{Argument}: The objects are genuinely distinct. CLIO studies
compression $\pi: X \to M$ and inference. MEM\textbar{}8 studies
event sequences $H_t$. Simulated Agency studies trajectory
optimization $\arg\max_a \mathbb{E}[\Omega]$. HYDRA studies
distributed module coordination. The analogy is probability theory,
statistics, decision theory, and control theory: shared substrate,
pedagogically distinct.

\textbf{Unifying formal object}: An agent is a tuple
$(\mathcal{A}, \pi, \mu, \mathcal{G})$ where $\mathcal{A}$ is its
admissibility space, $\pi$ its projection, $\mu$ its MEM\textbar{}8
history, and $\mathcal{G}$ its generative model. Each volume adds
one component to this tuple.

\textbf{Consequence}: Four-volume sequence (Volumes III--VII in
layer-centric curriculum). No merger.
\end{flybox}

\section{Candidate 7: Repair Scale Hierarchy}

\begin{flybox}
\textbf{Objects}: Ecphory (cognitive), Yarncrawler (semantic),
Institutional Repair (social), Xylomorphic Repair (physical).

\textbf{Verdict}: \textbf{SD} --- Shared Dynamical Form.

\textbf{Argument}: All are instances of the repair morphism
$R_\epsilon$ applied to different state spaces and domains. The
condition $d\Omega/dt > 0$ is the universal characterization at
all four scales.

\textbf{Scale table}:
\begin{center}
\begin{tabular}{lll}
Scale & Framework & Object repaired \\
\hline
Cognitive & Ecphory / MEM\textbar{}8 & Retrieval pathways \\
Semantic & Yarncrawler / Sem.\ Infra. & Knowledge structures \\
Institutional & Institutional Attractors & Institutional functions \\
Physical & Xylomorphic & Physical infrastructure \\
\end{tabular}
\end{center}

\textbf{Consequence}: Volume~VIII (Repair Theory) defines $R_\epsilon$
canonically. All other volumes cite this definition and specify their
domain-specific state space. This eliminates the central gap in
Volume~VIII while providing the connective tissue for Volumes V,
IX, and XII.
\end{flybox}

\section{Candidate 8: The Reachability Field}

\begin{flybox}
\textbf{Proposed object}: Reachability field $\mathfrak{R}(x)$
as a primitive, with derived quantities $\Omega(x) =
\mathrm{Vol}(\mathfrak{R}(x))$, $\mathcal{A}(x) = \log\Omega(x)$,
$\mathcal{I}(x) = -\log\Omega(x)$.

\textbf{Verdict}: \textbf{Promising but unconfirmed.}

\textbf{Argument for}: The identification $\mathcal{I}(x) =
-\log\Omega(x)$ unifies the tension landscape with the reachability
structure. If confirmed, repair ($d\Omega/dt > 0$), collapse
($d\Omega/dt < 0$), and agency ($\arg\max_a \mathbb{E}[\Omega]$)
all become instances of one geometric condition.

\textbf{Argument against (critical check)}: The RSVP entropy field
$S(x)$ encodes local accessibility density and the reachability
volume $\Omega(x)$ encodes global reachability. These are related
but not identical. Specifically:
$S(x) \sim -\nabla_x\log\Omega(x)$ is a gradient relationship
(accessibility gradient = entropy gradient), not an identity.
This means $\mathcal{I}(x) = -\log\Omega(x)$ must be verified
to be consistent with the RSVP field equations before adoption.

\textbf{Verification required}: Show that defining $\mathcal{I} =
-\log\Omega$ produces RSVP-compatible dynamics under the gradient
flow $\dot{x} = -\nabla\mathcal{I}$.

\textbf{Status}: Candidate identity (Claim~\ref{claim:ident}).
Present as conjecture until verified. Do not use as a theorem.
\end{flybox}

\section{Candidate 9: The Energy Functor}

\begin{flybox}
\textbf{Proposed object}: A functor $\mathfrak{T}:
\mathbf{Frameworks} \to \mathbf{EnergyLandscapes}$.

\textbf{Verdict}: \textbf{Philosophical proposal. Highest priority
for Volume XV.}

\textbf{Formal construction required}:

\noindent\textbf{Category $\mathbf{Frameworks}$}:
\begin{itemize}[nosep]
\item Objects: Flyxion frameworks $\mathcal{F} =
  (\mathcal{S}_\mathcal{F}, \mathcal{A}_\mathcal{F},
  g_\mathcal{F}, \mathcal{I}_\mathcal{F})$.
\item Morphisms: Admissibility-preserving smooth maps
  $\phi: \mathcal{S}_{\mathcal{F}_1} \to \mathcal{S}_{\mathcal{F}_2}$
  with $\phi(\mathcal{A}_{\mathcal{F}_1}) \subseteq
  \mathcal{A}_{\mathcal{F}_2}$.
\end{itemize}

\noindent\textbf{Category $\mathbf{EnergyLandscapes}$}:
\begin{itemize}[nosep]
\item Objects: Pairs $(\mathcal{S}, \mathcal{I})$.
\item Morphisms: Smooth $f: \mathcal{S}_1 \to \mathcal{S}_2$
  satisfying $\mathcal{I}_2(f(x)) \leq \mathcal{I}_1(x)$
  (inadmissibility non-increasing).
\end{itemize}

\noindent\textbf{Functor $\mathfrak{T}$}: $\mathfrak{T}(\mathcal{F})
= (\mathcal{S}_\mathcal{F}, \mathcal{I}_\mathcal{F})$; morphisms
map to the corresponding smooth maps on state spaces.

\textbf{What must be proved}:
(a) $\mathfrak{T}$ is well-defined (morphisms in
$\mathbf{Frameworks}$ map to morphisms in $\mathbf{EnergyLandscapes}$).
(b) $\mathfrak{T}$ is faithful (distinct frameworks give distinct
energy landscapes, i.e., the functor detects framework differences).
(c) $\mathfrak{T}$ is full (every admissibility-preserving map
between landscapes arises from a framework morphism).

\textbf{Consequence if confirmed}: Volume~XV is a theorem about
equivalence classes under $\mathfrak{T}$, not a narrative synthesis.
Frameworks $\mathcal{F}_1 \cong \mathcal{F}_2$ (same up to
morphism) iff $\mathfrak{T}(\mathcal{F}_1) \cong
\mathfrak{T}(\mathcal{F}_2)$ in $\mathbf{EnergyLandscapes}$.

\textbf{Status}: Define the categories formally. Prove (a) first.
(b) and (c) may require additional structure on $\mathbf{Frameworks}$.
\end{flybox}

% ================================================================
\chapter{The Master Claim and Path-Action Transition}
% ================================================================

\section{Two Levels of the Master Claim}

The corpus supports two formulations of the Master Claim at
different levels of generality.

\begin{masterbox}
\textbf{Master Claim, Level 1 (State-Functional Form)}:

For any Flyxion framework $\mathcal{F}$ with state space
$\mathcal{S}_\mathcal{F}$, admissibility space
$\mathcal{A}_\mathcal{F}$, metric $g_\mathcal{F}$, and
inadmissibility functional $\mathcal{I}_\mathcal{F}$, the
dynamics of $\mathcal{F}$ are governed by:
\[
\frac{dx}{dt} = -\nabla_{g_\mathcal{F}}
\mathcal{I}_\mathcal{F}(x) + \xi_t + \text{source terms},
\]
where $\xi_t$ is exploratory noise and source terms are
framework-specific.

\textbf{Status}: Confirmed as a structural organizing principle
by the formal spine document. Not yet a theorem because
``source terms'' are framework-specific and $\mathcal{I}_\mathcal{F}$
is derived from each framework rather than axiomatically imposed.
Becomes a theorem if the Energy Functor $\mathfrak{T}$ is proved
faithful.
\end{masterbox}

\begin{masterbox}
\textbf{Master Claim, Level 2 (Path-Variational Form)}:

For any Flyxion framework $\mathcal{F}$ whose dynamics are
path-dependent (Spherepop, MEM\textbar{}8, Agency, Repair,
Yarncrawler, TARTAN, HYDRA), the dynamics are governed by:
\[
\gamma^* = \operatorname{arg\,ext}_{\gamma \in \Gamma
(\mathcal{A}_\mathcal{F})} \mathcal{J}_\mathcal{F}[\gamma]
\]
where $\mathcal{J}_\mathcal{F}[\gamma] = \int_0^T
L_\mathcal{F}(\gamma,\dot\gamma,t)\,dt$ for a framework-specific
path Lagrangian $L_\mathcal{F}$.

\textbf{Status}: Structural claim. Level~2 subsumes Level~1 as
a special case.
\end{masterbox}

\section{Relationship Between the Two Levels}

\begin{proposition}[Level 1 as Special Case of Level 2]
Level~1 (gradient flow) is a special case of Level~2 (path
variational) when the path Lagrangian takes the form
$L(\gamma,\dot\gamma) = \frac{1}{2}g_{ij}\dot\gamma^i\dot\gamma^j
+ \mathcal{I}(\gamma)$
and boundary conditions $\gamma(0) = x_0$, $\gamma(T)$ free.
\end{proposition}
\begin{proof}
The Euler-Lagrange equations for this $L$ are:
$g_{ij}\ddot\gamma^j + \Gamma^i_{jk}\dot\gamma^j\dot\gamma^k
= -\nabla^i_g \mathcal{I}(\gamma)$,
which in the overdamped limit ($\ddot\gamma \to 0$) reduce to
$\dot\gamma^i = -g^{ij}\partial_j\mathcal{I}$, the gradient flow.
\end{proof}

\begin{warningbox}
\textbf{Critical asymmetry}: Level~2 does NOT reduce to Level~1 in
general. For path-action frameworks with irreversibility or
history-dependence (Spherepop, MEM\textbar{}8), the Lagrangian
encodes non-local path properties that cannot be expressed as
the gradient of any state functional. The proof strategy for
Volume~XV must use Euler-Lagrange arguments for these frameworks,
not gradient flow arguments. These are different proofs, not one
proof with two cases.
\end{warningbox}

\section{The Five-Level Hierarchy}

\begin{longtable}{p{0.5cm}p{3cm}p{3.5cm}p{5.5cm}}
\toprule
\textbf{L} & \textbf{Object} & \textbf{Symbol} & \textbf{Where frameworks diverge} \\
\midrule
0 & State spaces & $\mathcal{S}$ & All frameworks share this level \\
1 & Admissibility & $\mathcal{A} \subseteq \mathcal{S}$ & All frameworks share this level \\
2 & Trajectory spaces & $\Gamma(\mathcal{A})$ & \textbf{Branching point}. State-functional
  frameworks (RSVP, CLIO) do not require $\Gamma$. Path-action frameworks define their
  core objects here. \\
3 & Projection & $\pi: X \to M$ & Observation frameworks (CLIO, Agency) add this. \\
4 & Energy functionals & $\mathcal{I}$ & Different frameworks instantiate with different
  $\mathcal{I}$. Candidate~3 claims these are instances of one object. \\
5 & Dynamics & $\dot x = -\nabla_g \mathcal{I}$ & The master equation; special case of
  path variational. \\
\bottomrule
\end{longtable}

\section{Proof Obligations for Volume XV}

Volume~XV earns its place if and only if it proves that the
frameworks are related by explicit morphisms, not merely
analogous. The following proof obligations must all be discharged:

\medskip
\begin{enumerate}[leftmargin=1.5em]
\item \textbf{Construct $\mathbf{Frameworks}$ and
  $\mathbf{EnergyLandscapes}$} as formal categories (see
  Candidate~9 brief).

\item \textbf{Define $\mathfrak{T}$ explicitly} on objects and
  morphisms and verify it is a functor (preserves identity
  morphisms and composition).

\item \textbf{Prove $\mathfrak{T}$ is faithful}: if
  $\mathfrak{T}(\phi_1) = \mathfrak{T}(\phi_2)$ for morphisms
  $\phi_1, \phi_2: \mathcal{F}_1 \to \mathcal{F}_2$, then
  $\phi_1 = \phi_2$.

\item \textbf{Classify equivalence classes} under $\mathfrak{T}$:
  which frameworks are isomorphic in $\mathbf{EnergyLandscapes}$?

\item \textbf{State the synthesis theorem}:

\begin{center}
\fbox{\parbox{0.8\textwidth}{
Every persistent Flyxion system can be represented as an
admissibility-preserving transformation process. The category
of Flyxion frameworks embeds into $\mathbf{EnergyLandscapes}$
via the Energy Functor $\mathfrak{T}$, and frameworks
$\mathcal{F}_1 \cong \mathcal{F}_2$ iff
$\mathfrak{T}(\mathcal{F}_1) \cong \mathfrak{T}(\mathcal{F}_2)$.
}}
\end{center}

\item \textbf{Handle path-action frameworks}: the synthesis theorem
  must cover path-variational systems (Level~2 master claim), not
  only gradient flows (Level~1). This requires either extending
  $\mathbf{EnergyLandscapes}$ to include path Lagrangians as
  objects, or proving that every path-action framework has an
  equivalent state-functional reformulation (which is false in
  general, see §5.2).
\end{enumerate}

\section{What the Path-Action Transition Changes}

The recognition that the corpus contains both state-functional and
path-action frameworks has immediate consequences for Volume
organization:

\medskip
\begin{longtable}{p{4cm}lp{6.5cm}}
\toprule
\textbf{Framework} & \textbf{Form} & \textbf{Consequence for writing} \\
\midrule
RSVP & State & PDE-based; gradient flow language appropriate \\
CLIO & State & Fiber navigation; projection language appropriate \\
Spherepop & Path & Event irreversibility; Lagrangian
  $L = -\log p(e_t|H_t)$ \\
MEM\textbar{}8 & Path & History as primary object; Euler-Lagrange
  for viscous wave \\
Agency & Path & Trajectory selection;
  $\arg\max \mathbb{E}[\Omega(\gamma(T))]$ \\
Repair & Path & Trajectory correction; $L_\mathrm{rep}$ Lagrangian \\
Ecphory & Path & Wave propagation; activation threshold \\
Yarncrawler & Path & Route completion; divergence control \\
TARTAN & Path & Trajectory preservation; decomposition \\
Coordination & Mixed & Phase states; synchronization is path property \\
\bottomrule
\end{longtable}

% ================================================================
\chapter*{Closing Note: When Writing May Begin}
\addcontentsline{toc}{chapter}{Closing Note}
% ================================================================

Volume~0 is complete when the following have been addressed:

\begin{enumerate}
\item The symbol table in Chapter~1 has been adopted universally.
  All existing documents have been updated to use canonical symbols.

\item Every result flagged \textbf{As} in Chapter~3 has been
  either (a) provided with a proof or proof sketch, (b) downgraded
  to \textbf{Ph} (philosophical proposal), or (c) removed from
  the volumes that depend on it.

\item The foundational claims have been resolved as follows:
  the Reachability Principle (Claim~\ref{claim:rp}) is now
  \textbf{Sk} --- the contraction-rate exponential decay argument
  in Institutional Attractors Appendix~A constitutes a sketch,
  with two proof gaps remaining (necessity vs.\ sufficiency of
  $\Lambda \leq 0$, and universality across objective classes).
  The Reachability-Inadmissibility Identity (Claim~\ref{claim:ident})
  remains \textbf{As} and must be verified against the RSVP $S$ field
  before adoption. The Optimal Opacity claim is now \textbf{Sk} ---
  the optimum derivation is complete; the functional form derivation
  from first principles is the remaining obligation. The AI
  Convergence Certification (Claim~\ref{claim:ai-cert}) is newly
  added with status \textbf{As}; sufficiency of the three-condition
  conjunction must be proved before Volume~III is finalized.

\item The repair morphism (Definition~\ref{def:repair}) has been
  given a satisfying construction: either an explicit flow map or
  a proof that the repair Lagrangian Euler-Lagrange equations are
  well-posed.

\item The Equivalence Candidate briefs in Chapter~4 have been
  reviewed and their verdicts adopted as the canonical position.
  In particular: Candidate~1 is a common parent (not an
  equivalence); Candidates~3 and~5 are shared dynamical forms
  pending the Energy Functor; Candidate~9 is a philosophical
  proposal requiring formal construction.

\item The proof obligations for Volume~XV listed in §5.4 have
  been reviewed and a research plan for each has been assigned.
\end{enumerate}

\begin{warningbox}
None of the fifteen volumes should be submitted for publication
until Volume~0 is complete and all items above are resolved.
This is not a bureaucratic requirement. It is the condition under
which the fifteen volumes form a coherent field rather than a
library of theories.
\end{warningbox}

% ================================================================
\end{document}
% ================================================================
